% CICIL 2026 — Preliminary paper (review version). Compile with:
%   pdflatex cicil_prelim ; bibtex cicil_prelim ; pdflatex cicil_prelim ; pdflatex cicil_prelim
\pdfoutput=1
\documentclass[11pt]{article}
% Final (named) version: omit "review" so author names show and line numbers are
% off. Re-add [review] if the course wants line numbers (but that anonymizes authors).
\usepackage{ACL2023}
\usepackage{times}
\usepackage{latexsym}
\usepackage[T1]{fontenc}
\usepackage[utf8]{inputenc}
\usepackage{microtype}
\usepackage{inconsolata}
\usepackage{booktabs}
\usepackage{graphicx}

\title{Culturally-Aware Image Captioning for Indigenous Languages of the Americas:\\
A Two-Stage Vision-and-Retrieval Pipeline}

\author{Romit Basak \and Mehak Santosh Talmale \and Nandita Mahendra \and Tisha Patel \\
  Khoury College of Computer Sciences, Northeastern University \\
  \texttt{\{basak.r, talmale.m, mahendra.na, patel.tisha\}@northeastern.edu}}

\begin{document}
\maketitle

\begin{abstract}
Generic vision–language models (VLMs) discard culturally specific visual
information before it can reach a downstream translator, a bottleneck exposed
but left unresolved by the winning system of the AmericasNLP 2026 Cultural
Image Captioning for Indigenous Languages (CICIL) shared task. We describe a
two-stage pipeline that attacks this gap at the vision stage: Stage~1 uses a
structured \emph{cultural visual-question-answering} (VQA) prompt to elicit a
culturally grounded Spanish description together with per-category cultural
annotations; Stage~2 uses those annotations as retrieval keys and translates to
the target language with Gemini~2.5~Flash. We report preliminary results over
five languages (Guaraní, Bribri, Yucatec Maya, Wixárika, Nahuatl). On an
intrinsic Spanish-side proxy (Wixárika pilot, $n{=}20$) concise cultural-VQA
matches a generic baseline (20.9 vs.\ 21.0 ChrF++). End-to-end, cultural-VQA
yields no consistent improvement over the generic control: four of five
languages fall within $\pm 0.6$ ChrF++ (noise level) and only Nahuatl shows a
clear gain ($+1.56$), for a mean $\Delta$ of $+0.20$. We discuss why—including
an incomplete retrieval bank and greedy-decoding degeneration on the lowest-
resource languages—and outline the controlled experiments these preliminary
findings motivate.
\end{abstract}

\section{Introduction}
The CICIL shared task \citep{bui2026cicil} asks systems to caption culturally
situated images directly in an indigenous language of the Americas (Guaraní,
Yucatec Maya, Bribri, or Wixárika). The winning ``Gators'' submission
\citep{dhawan2026gators} established a strong two-stage recipe—generate a
Spanish intermediate description with a generic VLM, then translate with a
retrieval-augmented LLM—but its own analysis identifies the \emph{vision} stage
as the bottleneck: a generic Qwen2.5-VL encoder \citep{qwen2025qwen25vl}
produces the Spanish intermediate, so culturally specific semantics are stripped
out before translation begins.

We target exactly this gap. Our pipeline keeps the two-stage structure but makes
Stage~1 culturally aware through a structured cultural-VQA prompt that
interrogates the image (ceremony, material culture, landscape, kinship) before
synthesizing a Spanish description, and emits the per-category answers as
structured annotations that become Stage~2 retrieval keys. We ask three
research questions. \textbf{RQ1:} does culturally-aware vision encoding improve
end-to-end caption quality over a generic-VLM baseline? \textbf{RQ2:} does the
benefit vary with language resource level? \textbf{RQ3:} which cultural
categories are hardest to ground, and does this track caption quality? This
preliminary paper contributes (i) the pipeline and its cultural-VQA module,
(ii) an exploratory analysis of the CICIL data, and (iii) first end-to-end
baseline numbers—reported honestly, including where the intervention does
\emph{not} help.

\section{Related Work}

\subsection{Vision and Cultural Grounding}
Generic VLMs trained on Western, English-dominated web data carry systematic
cultural blind spots. Contrastive pretraining with CLIP \citep{radford2021clip}
yields cross-modal representations skewed toward that distribution, and
\citet{liu2021marvl} formalize the downstream cost through MaRVL, showing that
cross-lingual, cross-cultural visual reasoning falls far short of English
baselines. \citet{romero2024cvqa} confirm with CVQA—co-created with native
speakers across 30 countries—that culturally diverse visual QA remains open for
state-of-the-art models. This representational bottleneck is what
\citet{dhawan2026gators} uncover but leave unresolved.

Prior interventions typically act at the vision encoder.
\citet{huang2025cultureclip} build CulTwin, a synthetic set of visually similar
but culturally distinct triplets, and fine-tune CLIP with a tailored contrastive
objective for up to 5.49\% gains in fine-grained cultural concept recognition.
Contrastive encoder retraining, however, needs large curated image sets that our
$\sim$50-example-per-language budget cannot support; we therefore freeze the
SigLIP encoder \citep{zhai2023siglip} of SmolVLM
\citep{marafioti2025smolvlm} and inject cultural signal into the decoder via
prompting and LoRA \citep{hu2022lora}. A complementary line queries visual
content before captioning: the CIC framework \citep{yun2024cic} extracts
cultural elements via VQA and synthesizes a description with an LLM. Our
cultural-VQA module is a direct descendant, but specializes the question set to
the Americas and, crucially, repurposes its outputs as retrieval keys rather
than as terminal captions. \citet{li2026ravenea} (RAVENEA) show that
\emph{culture-aware retrieval} improves lightweight VLMs by $\geq$3.2\% on
cultural VQA and 6.2\% on cultural captioning, motivating our design in which
cultural grounding steers—rather than runs parallel to—retrieval. Finally,
category-stratified benchmarks (CulturalVQA, \citealp{nayak2024culturalvqa};
GIMMICK, \citealp{schneider2025gimmick}) report that tangible artifacts are
grounded more reliably than intangible practices, giving us a prior for RQ3 and,
with cultural theory \citep{yadav2025cultural}, a basis for our human-evaluation
rubric.

\subsection{Cultural Taxonomies}
\citet{hershcovich2022culture} frame culture in NLP along linguistic form,
values, common ground, and aboutness, arguing that serving a language means
serving its cultural context—captioning in Guaraní or Bribri is thus a cultural
grounding problem, not only translation. \citet{liu2025culturaltaxonomy} offer a
finer taxonomy grounded in anthropology and survey where NLP resources are
missing (non-Western, endangered, oral communities). We operationalize these
descriptive frameworks as a small set of \emph{visual question prompts}
(ceremony, material culture, landscape, kinship) derived, following CIC
\citep{yun2024cic}, from ethnographic documentation of the target communities.

\subsection{Retrieval, Prompting, and Evaluation}
\citet{agrawal2023icl} show that semantic similarity and diversity of in-context
examples both matter for MT, motivating retrieval by cultural concept rather than
surface text. \citet{merx2024mambai} obtain strong low-resource gains from
retrieval-augmented LLM prompting for Mambai but document a sharp in-domain vs.\
out-of-domain gap—directly relevant given that our indigenous corpora are drawn
from formal/religious registers \citep{mager2021americasnlp}. Our retrieval uses
multilingual Sentence-BERT embeddings \citep{reimers2020multilingual}, which
\citet{nie2022parc} find outperform sparse BM25 for morphologically rich
low-resource pairs. For evaluation we adopt ChrF++ \citep{popovic2017chrfpp}, the
official CICIL metric \citep{bui2026cicil}, which awards partial credit for
morphological variants; \citet{kumar2026chrf} caution that character metrics can
mask hallucination, source-copying, and repetition, which we corroborate below.
Our reporting layout follows per-language templates such as SemEval-2025 Task~11
\citep{muhammad2025semeval} and the AmericasNLP translation tasks
\citep{ebrahimi2023americasnlp,ebrahimi2024americasnlp}, and our human-evaluation
design follows the transparent, rubric-piloted protocol of THumB
\citep{kasai2022thumb}.

\section{Data and Exploratory Analysis}
The CICIL pilot/dev data are extremely small: the dev split provides 50 images
each for Guaraní, Bribri, Yucatec Maya, and Nahuatl, while Wixárika is supplied
only as a 20-image \emph{pilot}—and that pilot is the sole split carrying gold
Spanish captions (used as our Stage~1 proxy). Stage~2 retrieval banks come from
prior shared tasks and vary by orders of magnitude
(Table~\ref{tab:eda}): Guaraní is best resourced ($\sim$53k pairs,
\citealp{ebrahimi2023americasnlp}), Bribri and Wixárika far smaller
($\sim$7.5k/$\sim$9k, \citealp{mager2021americasnlp}), and Maya/Nahuatl banks are
still being assembled. Caption lengths differ substantially (Guaraní 21.6 words
on average vs.\ Maya 11.3); because the targets differ in morphological
typology, raw word counts are not comparable across languages, reinforcing our
use of character-level ChrF++ and cautioning against cross-language score
comparison.

\begin{table}[t]
\centering
\small
\begin{tabular}{lrrr}
\toprule
\textbf{Language} & \textbf{Dev} & \textbf{Cap.\ len} & \textbf{Bank pairs} \\
\midrule
Guaraní       & 50            & 21.6 & 53{,}183 \\
Bribri        & 50            & 14.9 & 7{,}506 \\
Yucatec Maya  & 50            & 11.3 & TBD \\
Nahuatl       & 50            & --   & TBD \\
Wixárika      & 20$^{\dagger}$ & 13.7 & 8{,}967 \\
\bottomrule
\end{tabular}
\caption{Exploratory analysis. ``Dev'' = examples in the split; ``Cap.\ len'' =
mean gold caption length (words); ``Bank pairs'' = Stage~2 retrieval corpus size
reported by the source shared task. $^{\dagger}$Wixárika ships as a 20-image
pilot (the only split with gold Spanish).}
\label{tab:eda}
\end{table}

\section{System}

\subsection{Stage 1: Culturally-Aware Vision}
Stage~1 turns an image into a culturally grounded Spanish description. The
\textbf{cultural-VQA} module first asks category-specific questions (ceremony,
material culture, landscape, kinship), then synthesizes the answers into a single
concise Spanish sentence and retains the per-category answers as structured
annotations. Our primary model is SmolVLM-2B \citep{marafioti2025smolvlm}: we
freeze its SigLIP encoder \citep{zhai2023siglip} and adapt only the language
decoder with LoRA \citep{hu2022lora} ($r{=}16$, $\alpha{=}32$, dropout $0.05$),
which is tractable on a single T4 within a \$50 budget; the frozen-encoder design
is exactly why cultural signal must enter through prompting rather than encoder
re-alignment. As the RQ1 control we run an off-the-shelf Qwen2.5-VL
\citep{qwen2025qwen25vl} with no cultural prompting. For these preliminary
numbers, both arms are produced with Qwen2.5-VL-7B served locally; the LoRA
fine-tune is in progress.

\subsection{Stage 2: Culturally-Indexed Retrieval and Translation}
Stage~2 encodes the Spanish side of each parallel corpus with
\texttt{paraphrase-multilingual-MiniLM-L12-v2}
\citep{reimers2020multilingual} into a FAISS inner-product index. The key design
choice is that queries are built from Stage~1's \emph{cultural annotations} (for
the generic arm, from the description itself), so retrieval targets cultural
rather than surface similarity \citep{agrawal2023icl,li2026ravenea}. The top-$k$
Spanish–indigenous pairs are supplied as few-shot demonstrations—most similar
last \citep{agrawal2023icl}—to Gemini~2.5~Flash, which translates the Stage~1
Spanish into the target language. No LLM fine-tuning is performed. We score with
sacreBLEU ChrF++ (\texttt{word\_order=2}) matching the official CICIL scorer.

\section{Preliminary Results}
\paragraph{Official baseline (reference).} Recomputed with the official scorer,
the shared-task baseline reaches 20.82 ChrF++ (Guaraní), 17.77 (Wixárika), 11.53
(Nahuatl), and 7.57 (Bribri); Maya has no MT baseline. These are the reference
points in Table~\ref{tab:rq1}.

\paragraph{Stage 1 intrinsic quality.} Because the authors do not read the target
languages, we first evaluate Stage~1 on its Spanish output against gold Spanish on
the Wixárika pilot ($n{=}20$), isolating vision-and-grounding from translation
(Table~\ref{tab:stage1}). A first cultural-VQA implementation
\emph{underperformed} the generic baseline (16.2 vs.\ 21.0), which diagnostic
analysis attributed to verbosity dilution—long, hedged syntheses (mean 105
tokens vs.\ $\sim$35) penalized by character-F against short references.
Constraining synthesis to one concise sentence (mean 25 tokens) closed the gap
(20.9 vs.\ 21.0, $\Delta{=}-0.15$; parity) while removing internally
inconsistent descriptions (0/20 vs.\ 2/20 naming incompatible materials). On this
proxy, concise cultural-VQA matches the generic baseline without degrading it.

\begin{table}[t]
\centering
\small
\begin{tabular}{lrrr}
\toprule
\textbf{Stage 1 configuration} & \textbf{ChrF++} & \textbf{len} & \textbf{conflicts} \\
\midrule
Generic baseline          & 21.0 & 35  & 2/20 \\
Cultural-VQA, verbose     & 16.2 & 105 & --   \\
Cultural-VQA, concise     & 20.9 & 25  & 0/20 \\
\bottomrule
\end{tabular}
\caption{Stage~1 intrinsic quality on the Wixárika pilot ($n{=}20$),
Qwen2.5-VL-7B, temperature 0, ChrF++ vs.\ gold Spanish. ``len'' = mean length
(tokens); ``conflicts'' = descriptions naming incompatible materials.}
\label{tab:stage1}
\end{table}

\paragraph{End-to-end (RQ1).} Table~\ref{tab:rq1} reports dev-set ChrF++ for the
full pipeline in both arms at $k{=}5$. \textbf{Cultural-VQA does not yield a
consistent end-to-end improvement.} Four of five languages fall within
$\pm0.6$ ChrF++ of the generic control—within run-to-run noise for $n{=}50$ and
therefore indistinguishable—and only Nahuatl shows a clear gain ($+1.56$); the
mean $\Delta$ across languages is $+0.20$. Two sanity checks pass: our generic
Guaraní (21.26) edges the official baseline (20.82) and generic Nahuatl (13.96)
exceeds its baseline (11.53), so the pipeline is behaving sensibly rather than
producing noise.

\begin{table}[t]
\centering
\small
\begin{tabular}{lrrrr}
\toprule
\textbf{Language} & \textbf{Base} & \textbf{Gen.} & \textbf{Cult.} & \textbf{$\Delta$} \\
\midrule
Guaraní       & 20.82 & 21.26 & 20.70 & $-0.56$ \\
Bribri        & 7.57  & 4.91  & 4.41  & $-0.50$ \\
Yucatec Maya  & --    & 19.25 & 19.56 & $+0.30$ \\
Wixárika      & 17.77 & 8.97  & 9.19  & $+0.22$ \\
Nahuatl       & 11.53 & 13.96 & 15.52 & $+1.56$ \\
\midrule
Mean $\Delta$ &       &       &       & $+0.20$ \\
\bottomrule
\end{tabular}
\caption{End-to-end dev ChrF++ ($k{=}5$). ``Base'' = official baseline;
``Gen.''/``Cult.'' = generic vs.\ cultural-VQA arms; $\Delta = $ Cult.\ $-$ Gen.
Four of five deltas are within noise ($\pm0.6$); only Nahuatl moves clearly.}
\label{tab:rq1}
\end{table}

\paragraph{Interpretation.} Three caveats bound these numbers. First, a genuine
retrieval bank exists \emph{only} for Wixárika (the 20-pair pilot); the other four
languages were translated zero-shot because their parallel banks are not yet
integrated, so the $k$-ablation is a single point ($k{=}5$) and the retrieval
contribution is untested for four of five languages. Second, greedy decoding
degenerated into repetition loops on the lowest-resource languages (Bribri
$\sim$5 ChrF++ in both arms; Nahuatl's \emph{generic} arm frequently looped),
which we bounded with an output-length cap—an honest RQ2 signal that a general
model cannot translate these languages without a real bank. Third, and
consequently, Nahuatl's $+1.56$ is confounded: cultural-VQA there partly wins by
avoiding a pathologically degenerate generic baseline rather than by demonstrably
better cultural grounding. Disentangling genuine grounding from reduced
degeneration is precisely what the planned human evaluation and category
breakdown (RQ3) are designed to do.

\section{Conclusion}
We presented a two-stage, culturally-aware captioning pipeline and its first
end-to-end numbers. On an intrinsic proxy, concise cultural-VQA matches a generic
baseline; end-to-end, it is presently indistinguishable from the control on four
of five languages, with one confounded gain. These results, together with an
incomplete retrieval bank, define the controlled experiments we prioritize next:
integrating the full parallel banks, running the LoRA fine-tune, and the
category-stratified and human evaluations that address RQ2 and RQ3.

\section*{Limitations}
Our targets are morphologically rich, low-resource languages; results are on
small dev/pilot splits ($n{\leq}50$), so ChrF++ differences below roughly one
point are not meaningful and we report them as such. The cultural-VQA arm's
end-to-end contribution is confounded by (i) a retrieval bank present for only
one language, (ii) reliance on a proprietary translator (Gemini~2.5~Flash) whose
behavior may change, and (iii) greedy-decoding degeneration on the lowest-resource
languages. The Stage~1 intrinsic evaluation exists only for Wixárika (the sole
split with gold Spanish), and the SmolVLM LoRA fine-tune is not yet complete, so
the reported arms use an off-the-shelf VLM. Several teammate-supplied citations
are transcribed from draft reference lists and marked for verification.

\section*{Ethics Statement}
The CICIL data are released under CC~BY-NC~4.0. To respect the NonCommercial
clause we translate via Vertex~AI, which does not use inputs to train models,
rather than a consumer API tier that may; our use is non-commercial academic
research and we attribute the dataset. This work aims to support, not supplant,
indigenous-language communities; captions generated by our system are not
authoritative and should be validated by native and heritage speakers before any
community-facing use. We report negative and confounded results transparently to
avoid overstating the utility of automated cultural captioning.

\bibliography{custom}
\bibliographystyle{acl_natbib}

\end{document}
